\pdfbookmark[0]{Section V: Quadratic forms and indeterminate equations}{section-five}
\SourceLeaf{125}
\begin{center}
SECTION FIVE

QUADRATIC FORMS AND

INDETERMINATE EQUATIONS OF THE SECOND DEGREE
\end{center}
\begin{center}\rule{.35\linewidth}{.4pt}\end{center}

\begin{center}
\textit{Aim of the investigation; definition and notation of forms.}
\end{center}

\Article{153}
In this section we will mainly study functions of two indeterminates $x,y$ of the form
\[
axx+2bxy+cyy
\]
where $a,b,c$ are given integers. We will call these \textit{forms of the second degree}, or simply \textit{forms}. This investigation will provide the basis for solving the famous problem of finding all solutions of any indeterminate equation of the second degree in two unknowns, whether the unknowns are required to have integer values or only rational values. The distinguished Lagrange has already solved this problem in full generality. Many further results about the nature of \textit{forms} have been discovered for the first time, either by this great geometer himself or by the distinguished Euler; others, discovered earlier by Fermat, have been supplied with proofs by them. But as we pursued our investigation of forms closely, so many new results presented themselves that we thought it worth taking up the whole subject again from the beginning. We had all the more reason to do so because we found that the discoveries of these men, scattered over many places, were known to few people; because the method we will use is largely our own; and finally
\SourceLeaf{126}
because our results could not even be understood without a fresh account of theirs. We have no doubt that many excellent results in this field still lie hidden, on which others can try their strength. We will always give the history of notable results in its proper place.

We will denote the form $axx+2bxy+cyy$, when the indeterminates $x,y$ are not at issue, by $(a,b,c)$. This notation will thus stand, without specifying the indeterminates, for the sum of three parts: the product of the given number $a$ and the square of any indeterminate; the product of twice the number $b$, this indeterminate and another indeterminate; and the product of the number $c$ and the square of the second indeterminate. \textit{For example}, $(1,0,2)$ will denote the sum of a square and twice a square. Although the forms $(a,b,c)$ and $(c,b,a)$ denote the same thing if we consider only the \textit{parts themselves}, they differ if we also consider the \textit{order} of the parts. We will therefore distinguish them carefully from now on; the advantages of doing this will become clear enough in what follows.

\Needspace{8\baselineskip}
\begin{center}
\textit{Representation of numbers; determinant.}
\end{center}

\Article{154}
We will say that a given number is \textit{represented} by a given form if integer values are assigned to the indeterminates of the form so that its value equals the given number. We have the following

\textsc{Theorem.} \textit{If the number $M$ can be represented by the form $(a,b,c)$ so that the values of the indeterminates that give this representation are relatively prime, then $bb-ac$ is a quadratic residue of the number $M$.}

\textit{Proof.} Let the values of the indeterminates be $m,n$, so that
\[
amm+2bmn+cnn=M
\]
and choose numbers $\mu,\nu$ so that $\mu m+\nu n=1$ (Article 40). Expansion easily shows that
\[
\begin{aligned}
&amm+2bmn+cnn)(a\nu\nu-2b\mu\nu+c\mu\mu\\
&\qquad={}(\mu(mb+nc)-\nu(ma+nb))^2-(bb-ac)(m\mu+n\nu)^2
\end{aligned}
\]
or
\[
M(a\nu\nu-2b\mu\nu+c\mu\mu)=(\mu(mb+nc)-\nu(ma+nb))^2-(bb-ac)
\]
So
\SourceLeaf{127}
\[
bb-ac\equiv(\mu(mb+nc)-\nu(ma+nb))^2\ (\mathrm{mod.}\,M)
\]
that is, $bb-ac$ is a quadratic residue of $M$.

We will call the number $bb-ac$ the \textit{determinant} of the form $(a,b,c)$; we will show later that the properties of a form depend above all on the nature of this number.

\begin{center}
\textit{Values of the expression $\surd(bb-ac)(\mathit{mod.}\,M)$ to which a representation of the number $M$ by the form $(a,b,c)$ belongs.}
\end{center}

\Article{155}
Thus
\[
\mu(mb+nc)-\nu(ma+nb)
\]
is a value of the expression
\[
\surd(bb-ac)\ (\mathrm{mod.}\,M)
\]
Now the numbers $\mu,\nu$ can be chosen in infinitely many ways so that $\mu m+\nu n=1$, giving one value after another of that expression. Let us see how these values are related. Suppose not only that $\mu m+\nu n=1$, but also that $\mu'm+\nu'n=1$, and put
\[
\mu(mb+nc)-\nu(ma+nb)=v,\quad
\mu'(mb+nc)-\nu'(ma+nb)=v'
\]
Multiplying the equation $\mu m+\nu n=1$ by $\mu'$, the other equation $\mu'm+\nu'n=1$ by $\mu$, and subtracting gives $\mu'-\mu=n(\mu'\nu-\mu\nu')$. Similarly, multiplying the first by $\nu'$ and the second by $\nu$, and then subtracting, gives $\nu'-\nu=m(\mu\nu'-\mu'\nu)$. This immediately gives
\[
v'-v=(\mu'\nu-\mu\nu')(amm+2bmn+cnn)=(\mu'\nu-\mu\nu')M
\]
or $v'\equiv v\ (\mathrm{mod.}\,M)$. So however $\mu,\nu$ are chosen, the formula $\mu(mb+nc)-\nu(ma+nb)$ cannot give \textit{different} (that is, incongruent) values of the expression $\surd(bb-ac)\ (\mathrm{mod.}\,M)$. If $v$ is any value of that formula, we will say that the representation of the number $M$ by the form $axx+2bxy+cyy$ in which $x=m$, $y=n$ \textit{belongs to the value $v$ of the expression $\surd(bb-ac)(\mathit{mod.}\,M)$}. It is also easy to show that if one value of this formula is $v$ and $v'\equiv v\ (\mathrm{mod.}\,M)$, the numbers $\mu,\nu$ that give $v$ can be replaced by other numbers $\mu',\nu'$ that give $v'$. In fact, putting
\[
\mu'=\mu+\frac{n(v'-v)}{M},\quad
\nu'=\nu-\frac{m(v'-v)}{M}
\]
gives
\[
\mu'm+\nu'n=\mu m+\nu n=1
\]
\SourceLeaf{128}
and the value of the formula obtained from $\mu',\nu'$ exceeds the value obtained from $\mu,\nu$ by the quantity $(\mu'\nu-\mu\nu')M$. This quantity becomes $=(\mu m+\nu n)(v'-v)=v'-v$, giving the value $=v'$.

\Article{156}
If we have two representations of the same number $M$ by the same form $(a,b,c)$, in each of which the values of the indeterminates are relatively prime, they may belong either to the same value of the expression $\surd(bb-ac)\ (\mathrm{mod.}\,M)$ or to different values. Let
\[
M=amm+2bmn+cnn=am'm'+2bm'n'+cn'n'
\]
and
\[
\mu m+\nu n=1,\quad \mu'm'+\nu'n'=1
\]
Clearly, if
\[
\mu(mb+nc)-\nu(ma+nb)\equiv\mu'(m'b+n'c)-\nu'(m'a+n'b)\ (\mathrm{mod.}\,M)
\]
the congruence remains true however we replace $\mu,\nu$ and $\mu',\nu'$ by other suitable values. In this case we will say that the two representations belong to the \textit{same} value of the expression $\surd(bb-ac)\ (\mathrm{mod.}\,M)$. But if the congruence fails for any choice of $\mu,\nu$ and $\mu',\nu'$, it will hold for none, and the representations will belong to \textit{different} values. If, on the other hand,
\[
\mu(mb+nc)-\nu(ma+nb)\equiv-(\mu'(m'b+n'c)-\nu'(m'a+n'b))
\]
we will say that the representations belong to \textit{opposite} values of the expression $\surd(bb-ac)$. We will also use all these terms when considering several representations of the same number by \textit{different} forms with the same determinant.

\textit{Example.} Take the form $(3,7,-8)$, with determinant $=73$. It gives the following representations of the number $57$:
\[
3\mathbin{.}13^2+14\mathbin{.}13\mathbin{.}25-8\mathbin{.}25^2;\quad
3\mathbin{.}5^2+14\mathbin{.}5\mathbin{.}9-8\mathbin{.}9^2
\]
For the first we can put $\mu=2$, $\nu=-1$, giving the value of the expression $\surd73\ (\mathrm{mod.}\,57)$ to which the representation belongs:
\[
{}=2(13\mathbin{.}7-25\mathbin{.}8)+(13\mathbin{.}3+25\mathbin{.}7)=-4
\]
\SourceLeaf{129}
Similarly, putting $\mu=2$, $\nu=-1$, we find that the second representation belongs to the value $+4$. So the two representations belong to opposite values.

Before going further, we note that forms with determinant $=0$ are excluded entirely from the investigations that follow. They would only disrupt the orderly form of the theorems and so need a treatment of their own.

\begin{center}
\textit{A form containing another, or contained in another; proper and improper transformation.}
\end{center}

\Article{157}
If a form $F$, with indeterminates $x,y$, can be changed into another form $F'$, with indeterminates $x',y'$, by substitutions
\[
x=\alpha x'+\beta y',\quad y=\gamma x'+\delta y'
\]
in which $\alpha,\beta,\gamma,\delta$ are integers, we will say that the first \textit{contains} the second, or that the second \textit{is contained in the first}. Let the form $F$ be
\[
axx+2bxy+cyy
\]
and the form $F'$ be
\[
a'x'x'+2b'x'y'+c'y'y'
\]
We then have the following three equations:
\[
\begin{aligned}
a'&=a\alpha\alpha+2b\alpha\gamma+c\gamma\gamma\\
b'&=a\alpha\beta+b(\alpha\delta+\beta\gamma)+c\gamma\delta\\
c'&=a\beta\beta+2b\beta\delta+c\delta\delta
\end{aligned}
\]
Multiplying the second equation by itself and the first by the third, then subtracting and removing the terms that cancel, gives
\[
b'b'-a'c'=(bb-ac)(\alpha\delta-\beta\gamma)^2
\]
It follows that the determinant of the form $F'$ is divisible by the determinant of the form $F$, and that the quotient is a square. These determinants therefore have \textit{the same sign}. If, in addition, the form $F'$ can be changed into the form $F$ by a similar substitution, that is, if both $F'$ is contained in $F$ and $F$ in $F'$,
\SourceLeaf{130}
the determinants of the forms are equal\HistoricalNote{*)}{The preceding analysis clearly shows that this proposition also extends to forms with determinant $=0$. But the equation $(\alpha\delta-\beta\gamma)^2=1$ must not be extended to that case.} and $(\alpha\delta-\beta\gamma)^2=1$. In this case we will call the forms \textit{equivalent}. Equality of determinants is therefore a necessary condition for equivalence of forms, although equivalence certainly does not follow from it alone. --- We will call the substitution $x=\alpha x'+\beta y'$, $y=\gamma x'+\delta y'$ a \textit{proper transformation} if $\alpha\delta-\beta\gamma$ is positive, and an \textit{improper transformation} if $\alpha\delta-\beta\gamma$ is negative. We will say that the form $F'$ is \textit{properly} or \textit{improperly} contained in the form $F$ if $F$ can be changed into $F'$ by a proper or improper transformation. Thus, if the forms $F,F'$ are equivalent, $(\alpha\delta-\beta\gamma)^2=1$, so that for a proper transformation $\alpha\delta-\beta\gamma=+1$, and for an improper one $=-1$. --- We will call several transformations \textit{similar} if they are all proper or all improper; a proper and an improper transformation will be called \textit{dissimilar}.

\begin{center}
\textit{Proper and improper equivalence.}
\end{center}

\Article{158}
\textit{If the forms $F,F'$ have equal determinants and $F'$ is contained in $F$, then $F$ is also contained in $F'$, properly or improperly according as $F'$ is properly or improperly contained in $F$.} Suppose $F$ changes into $F'$ on putting
\[
x=\alpha x'+\beta y',\quad y=\gamma x'+\delta y'
\]
Then $F'$ changes into $F$ on putting
\[
x'=\delta x-\beta y,\quad y'=-\gamma x+\alpha y
\]
For this substitution clearly produces from $F'$ the same result as putting in $F$
\[
x=\alpha(\delta x-\beta y)+\beta(-\gamma x+\alpha y),\quad
y=\gamma(\delta x-\beta y)+\delta(-\gamma x+\alpha y)
\]
or
\[
x=(\alpha\delta-\beta\gamma)x,\quad y=(\alpha\delta-\beta\gamma)y
\]
This clearly changes $F$ into $(\alpha\delta-\beta\gamma)^2F$, that is, back into $F$ (the preceding article). The second transformation is clearly proper or improper according as the first is proper or improper.

If both $F'$ is \textit{properly} contained in $F$ and $F$ in $F'$, we will call the forms \textit{properly \SourceLeaf{131}equivalent}; if each is improperly contained in the other, we will call them \textit{improperly equivalent}. --- The use of these distinctions will soon become clear.

\textit{Example.} The form $2xx-8xy+3yy$ changes, under the substitutions $x=2x'+y'$, $y=3x'+2y'$, into the form $-13x'x'-12x'y'-2y'y'$, and the latter changes into the former on putting $x'=2x-y$, $y'=-3x+2y$. So the forms $(2,-4,3)$, $(-13,-6,-2)$ are \textit{properly equivalent}.

\Needspace{12\baselineskip}
The problems we will now take up are these. I. Given any two forms with the same determinant, determine whether they are equivalent, and whether properly, improperly, or in both ways, since this too is possible. If their determinants are unequal, determine whether at least one contains the other, properly, improperly, or in both ways. Finally, find all transformations of one into the other, both proper and improper. II. Given any form, determine whether a given number can be represented by it, and find all its representations. But forms with negative determinant need a different method here from forms with positive determinant. We will therefore first give what is common to both, then consider forms of each kind separately.
